\subsection{Mechanical Tests}
\subsubsection{Gantry Tilt: Accuracy of the tilt Indicator}

Most CT scanners accomplish non-orthogonal scans by tilting the gantry with some angle. If the tilt is not accurate, the gantry may hit the patient for some extreme clinical conditions.
For the measurement radiochromic film was used.
\begin{figure}[H]
    \centering
    \includegraphics[width=0.3\textwidth]{/Users/souvik/Desktop/Lab-Reports/4th Sem/Quality Assurance of a Diagnostic CT Machine/Figures/Experimental/gantrytilt.jpeg}
    \caption{Gantry Tilt Measurement using Radiochromic Film}
    \label{fig:gantry_tilt}
\end{figure}
Actual gantry tilt = ±15 degrees
Measured gantry tilt = 15.8, 15.2 degrees
Difference = 0.8, 0.2 degrees
Tolerance: The measured gantry tilt should be within ± 2 degree of the actual gantry tilt.
So, the measured gantry tilt is within the tolerance limit of ± 2 degree.

\subsection{Radiation Safety Test of Computed Tomography Equip-
ment}
\subsubsection{Slice Thickness}
Slice thickness in the Pro-CT Mini is determined by measuring the length of ramps in the
geometrical module image. The measured length depends on the slice thickness and an angle
at which the ramp is inclined.
\begin{figure}[H]
    \centering
    \includegraphics[width=0.6\textwidth]{/Users/souvik/Desktop/Lab-Reports/4th Sem/Quality Assurance of a Diagnostic CT Machine/Figures/Theory/Slice thickness.png}
    \caption{Geometry used in the Pro-CT Mini phantom for measurement of slice thickness}
    \label{fig:slice_thickness}
\end{figure}
In the pro-CT mini phantom, ramps are at 45 degree angle so $\tan(45^o) = 1$, hence the measured length of the ramp is equal to the slice thickness.It is taken from this \href{https://www.scribd.com/document/625663053/C-T-Slice-Thickness}{\textbf{ref}}.

\begin{figure}[H]
    \centering
    \includegraphics[width=0.6\textwidth]{/Users/souvik/Desktop/Lab-Reports/4th Sem/Quality Assurance of a Diagnostic CT Machine/Figures/Experimental/slicethickness.jpeg}
    \caption{Slice Thickness Measurement using Pro-CT Mini Phantom}
    \label{fig:slice_thickness}
\end{figure}

Set slice thickness = 5 mm
Measured slice thickness = 4.3 mm, 4.2 mm, 4.3 mm, 4.5 mm
Average measured slice thickness = 4.325 mm

AERB Tolerance:

For less than 1 mm : 0.5 

1 mm to 2 mm : ± 50 \%

Above 2 mm : ± 1 mm

So it is under the tolerance limit of ± 1 mm.

\subsubsection{Accuracy of Operating Potential}
\begin{figure}[H]
    \centering
    \includegraphics[width=0.6\textwidth]{/Users/souvik/Desktop/Lab-Reports/4th Sem/Quality Assurance of a Diagnostic CT Machine/Figures/Experimental/kVp accuracy.png}
\end{figure}

\subsubsection{Accuracy of Timer}
\begin{figure}[H]
    \centering
    \includegraphics[width=0.6\textwidth]{/Users/souvik/Desktop/Lab-Reports/4th Sem/Quality Assurance of a Diagnostic CT Machine/Figures/Experimental/TimerAccuracy.png}
\end{figure}
\subsubsection{Linearity of Radiation Output}
Operating potential = 120 kVp
\begin{figure}[H]
    \centering
    \includegraphics[width=0.6\textwidth]{/Users/souvik/Desktop/Lab-Reports/4th Sem/Quality Assurance of a Diagnostic CT Machine/Figures/Experimental/Linearity.png}
\end{figure}
\subsubsection{Reproducibility of Radiation Output}
Operating current = 100 mA
\begin{figure}[H]
    \centering
    \includegraphics[width=0.6\textwidth]{/Users/souvik/Desktop/Lab-Reports/4th Sem/Quality Assurance of a Diagnostic CT Machine/Figures/Experimental/Reproducibilit.png}
\end{figure}

\subsubsection{ Radiation Dose Test Measurement of Weighted CT Dose Index
($CTDI_w$)}
Operating potential = 120 kVp
Tickness = 0.5 cm
Operating current = 200 mA
\begin{figure}[H]
    \centering
    \includegraphics[width=0.6\textwidth]{/Users/souvik/Desktop/Lab-Reports/4th Sem/Quality Assurance of a Diagnostic CT Machine/Figures/Experimental/CTDI.png}
\end{figure}
Tolerance:
±20 \% of the quoted value for head phantom and ± 40 \% of the quoted value for body phantom.
\subsubsection{Low Contrast Resolution}
Phantom = Pro-CT Mini Phantom

window width = 395, window level = 210

Low Contrast Resolution = 2 mm at 0.5 \% contrast level.

\textbf{Tolerance:}
2.5 mm at 0.5 \% contrast difference (expected) or 5 mm at 1 \% contrast difference (minimum).
\begin{figure}[H]
    \centering
    \includegraphics[width=0.6\textwidth]{/Users/souvik/Desktop/Lab-Reports/4th Sem/Quality Assurance of a Diagnostic CT Machine/Figures/Experimental/Low contrast.jpeg}
    \caption{Low Contrast Resolution Measurement using Pro-CT Mini Phantom}
    \label{fig:low_contrast_resolution}
\end{figure}

The low contrast resolution is 2 mm at 0.5 \% contrast level which is within the tolerance limit
of 2.5 mm at 0.5 \% contrast difference, hence the low contrast resolution of the CT machine is
acceptable.
\subsubsection{High Contrast Resolution}
Phantom = Pro-CT Mini Phantom

window width = 490, window level = 307

Size of the smallest resolvable bar pattern = 5 lp/cm.

\textbf{Tolerance:}
At 10 \% contrast difference the size of the bar pattern that could be resolved should be 1.6 mm (3.12 lp/cm), expected high contrast resolution: 0.8 mm ($\approx$ 6.25 lp/cm).

\begin{figure}[H]
    \centering
    \includegraphics[width=0.6\textwidth]{/Users/souvik/Desktop/Lab-Reports/4th Sem/Quality Assurance of a Diagnostic CT Machine/Figures/Experimental/High constrast.jpeg}
    \caption{High Contrast Resolution Measurement using Pro-CT Mini Phantom. Max. resolved line pair pattern for measurement of high contrast resolution is 5 lp/cm.}
    \label{fig:low_contrast_resolution}
\end{figure}
\subsubsection{Radiation Leakage Level from X-ray Tube Housing}
Operating potential = 140 kVp
Operating current = 500 mA
time = 1 sec
\begin{figure}[H]
    \centering
    \includegraphics[width=0.6\textwidth]{/Users/souvik/Desktop/Lab-Reports/4th Sem/Quality Assurance of a Diagnostic CT Machine/Figures/Experimental/Leakage.png}
\end{figure}
\subsubsection{Radiation Protection Survey of the Installation}
\begin{figure}[H]
    \centering
    \includegraphics[width=0.6\textwidth]{/Users/souvik/Desktop/Lab-Reports/4th Sem/Quality Assurance of a Diagnostic CT Machine/Figures/Experimental/Survey.png}
\end{figure}
